\documentclass[11pt,a4paper]{letter}
\usepackage[margin=2.5cm]{geometry}
\usepackage[utf8]{inputenc}
\usepackage[T1]{fontenc}
\usepackage{hyperref}
\setlength{\parskip}{0.6\baselineskip}
\setlength{\parindent}{0pt}

\signature{Miguel Ángel Díaz-Campos\\Instituto Nacional de Pediatría\\Ciudad de México, México\\\href{mailto:mdiazc161@unam.edu}{mdiazc161@unam.edu}}
\address{Miguel Ángel Díaz-Campos\\Instituto Nacional de Pediatría (INP)\\Av.\ Insurgentes Sur 3700, Letra C\\04530 Coyoacán, Ciudad de México\\México}
\date{}

\begin{document}

\begin{letter}{Oficina Editorial\\\textit{Nature Communications}\\Springer Nature}
\opening{Estimado(a) Editor(a):}

Nos complace someter nuestro manuscrito titulado \textbf{«Análisis de consenso multi-método de transcriptómica espacial \textit{in situ} dirigida: un marco con demostración piloto sobre adenocarcinoma humano de pulmón»} para su consideración como Research Article en \textit{Nature Communications}.

\textbf{Qué hace el trabajo.} Presentamos un marco de consenso multi-herramienta para analizar datos de transcriptómica espacial \textit{in situ} dirigida y lo demostramos en un dataset Xenium público de adenocarcinoma pulmonar (268{,}034 células, panel de 289 genes). Para cada familia analítica --- anotación de tipos celulares, detección de genes espacialmente variables, inferencia de dominios espaciales, comunicación célula--célula y pseudotiempo --- ejecutamos al menos dos herramientas computacionales independientes y retenemos únicamente los hallazgos concordantes entre métodos. El marco rinde 288 de 289 genes del panel como genes espacialmente variables de consenso (prueba de tres métodos); un subconjunto de consenso por célula entre las particiones de dominio espacial de Banksy y Novae (ARI~$=0.15$ a nivel de partición, 35.4\% de células de alta confianza a nivel por célula); 20 de 30 interacciones célula--célula retenidas bajo Benjamini--Hochberg conjunto (LIANA+) y corrección de Bonferroni (Spacia); y un pseudotiempo de consenso de dos métodos (Slingshot + diffusion pseudotime, $r = 0.478$) restringido al compartimento macrofágico.

\textbf{Por qué importa.} El hub dominante validado bajo Bonferroni en comunicación célula--célula presenta a ADAM17 (sheddasa) y MUC1 (oncomucina), co-localizados espacialmente a través de seis tipos celulares emisores mieloides y estromales independientes convergentes sobre células tumorales, con $p$-values ajustados de Spacia que alcanzan $1.10 \times 10^{-71}$. La actividad tanto en estados macrofágicos M1 como M2, combinada con un continuo de activación macrofágico no direccional (ambos marcadores M1 y M2 aumentan con el pseudotiempo), apoya un modelo de inmunosupresión mediada por shedding, no por checkpoint, en la interfaz tumor--parénquima --- distinto del paradigma canónico de agotamiento por checkpoint y terapéuticamente accionable mediante inhibición de ADAM17.

\textbf{Por qué \textit{Nature Communications}.} El trabajo combina un marco metodológico con un estudio biológico de caso sustantivo. El marco es genérico a través de plataformas espaciales in situ y tamaños de panel, y documentamos explícitamente los límites metodológicos (qué herramientas fallan a 289 genes, por qué, y a qué tamaño de panel se espera que cada una opere; Nota Suplementaria S1). El pipeline open-source que acompaña está diseñado para portabilidad a la cohorte multi-muestra multi-órgano que nosotros y otros analizaremos en la siguiente ronda de estudios Xenium. El alcance metodológico combinado y la profundidad biológica encajan mejor con el lectorado cross-disciplinario de \textit{Nature Communications} que con un venue puramente metodológico.

\textbf{Transparencia.} Una auditoría tipo revisor identificó cinco problemas en un encuadre anterior de este trabajo (Nota Suplementaria S1.3): una prueba downstream aplicada a embeddings convertidos a counts enteros, una configuración de trayectoria implausible, un umbral de magnitud-rank etiquetado erróneamente como FDR, $p$-values per-test no corregidos sobre las pruebas cruzadas, y deriva de versiones de software entre código y manuscrito. Los cinco fueron resueltos antes del envío, y el memo de auditoría se incluye completo como nota suplementaria --- tanto como ayuda al revisor como una declaración de transparencia que, en nuestra opinión, el campo actualmente carece.

\textbf{Limitaciones honestas.} El piloto es una sola sección de un solo paciente. La significancia estadística a nivel celular es alta pero no sustituye a la replicación biológica. Enmarcamos este artículo como metodología + estudio de caso, con hallazgos biológicos generadores de hipótesis. Un estudio de cohorte multi-muestra usando el mismo marco sobre un panel de 439 genes (fibrosis pulmonar y hepática asociadas a disqueratosis congénita) está actualmente en recolección y se reportará por separado.

\textbf{Disponibilidad de datos y código.} El dataset está públicamente disponible en 10x Genomics (CG000709 Rev~C). El pipeline completo (DAG de Snakemake, entornos conda, imagen Docker para Spacia) está disponible en \url{https://github.com/MiguelDiaz02/xenium-spatial-transcriptomics} bajo licencia MIT; la versión definitiva se etiquetará en la aceptación.

El manuscrito no ha sido sometido ni considerado por ninguna otra revista. Todos los autores han aprobado el envío y declaran no tener conflictos de interés.

Agradecemos su consideración y quedamos atentos a su decisión.

\closing{Atentamente,}

\end{letter}

\end{document}
